\chapter*{General Conclusion and Perspectives}
\addcontentsline{toc}{chapter}{General Conclusion and Perspectives}
\label{chap:general_conclusion}

This graduation project focused on the design and implementation of \textbf{VisionTest-AI}, an intelligent functional testing agent that combines natural language understanding, browser automation, visual perception, and automatic reporting. The main objective was to reduce the fragility of traditional automated tests by allowing the system to understand Gherkin scenarios, generate executable actions, execute them with Playwright, and use visual fallback mechanisms when classical selectors fail.

Throughout the project, we built a complete prototype composed of several modules. The NLP module parses Gherkin scenarios, classifies intents, extracts entities, and prepares structured action payloads. The semantic mapping module transforms these actions into Playwright-compatible commands and selector candidates. The execution service runs the test steps in a browser, while the vision module uses YOLOv8n and OCR to detect UI elements from screenshots when selector-based execution is not sufficient. Finally, the reporting service generates JSON and HTML reports containing execution details, screenshots, errors, and summary metrics.

One of the main achievements of this work is the integration of several technologies into a single functional pipeline. The project does not only generate scripts or analyze screenshots separately; it connects these components into a complete workflow, starting from a business-readable Gherkin scenario and ending with an execution report. This demonstrates the feasibility of an intelligent testing assistant that can support QA engineers in writing, executing, and analyzing functional tests.

The project also highlighted several technical challenges. Preparing a reliable UI detection dataset was more difficult than expected. Ready-made datasets were not diverse enough, manual annotation was time-consuming, and automatic labeling required several corrections before becoming usable. The final DOM-assisted labeling approach helped improve the quality of annotations, but the YOLO model still needs more training data and refinement. Deployment was another challenge, especially because the system requires browser automation, OCR, AI models, and report storage, which makes it heavier than a simple Streamlit application.

The obtained results are encouraging. The NLP and semantic mapping modules successfully handled explicit functional steps such as navigation, input, click, and verification. The system was able to generate Playwright scripts, execute scenarios, capture screenshots, and produce structured JSON and HTML reports. The visual fallback mechanism also showed its value in some cases, especially when the target element was visible but the selector-based strategy failed.

However, the current version remains a prototype and has several limitations. The vision module is not yet fully reliable on all interfaces, OCR can fail on small or low-contrast text, and some complex interactions such as drag-and-drop, file upload, or advanced dynamic components are not fully supported. The system also does not yet include production features such as authentication, project history, PostgreSQL storage, Redis queues, or advanced cloud orchestration.

Several perspectives can be considered for future work. The NLP module can be improved by adding multilingual support, especially for French and mixed-language Gherkin scenarios. The vision dataset can be expanded with more annotated examples from different websites, themes, and component libraries. The execution engine can be enhanced with better ranking logic, stronger self-healing strategies, and support for more complex user actions. On the platform side, adding PostgreSQL, Redis, authentication, and project management features would make the system closer to an industrial solution.

In conclusion, \textbf{VisionTest-AI} provides a solid foundation for an AI-assisted functional testing platform. It shows that combining semantic understanding with visual perception can make test automation more flexible and more resilient. Although further improvements are still required, the project successfully demonstrates how AI techniques can be applied to modern QA workflows in a practical and extensible way.